% from ega4-1-fr.tex
\oldpage[IV-1]{250}
où $t$ parcourt $B-\mathfrak q$. On a, par suite (1.9.2),
\[
  f(\operatorname{Spec}(\mathcal O_{X,x}))
  =\bigcap_t f(\operatorname{Spec}(B_t))=\bigcap_t f(\mathrm D(t)),
\]
et par hypothèse, pour tout $t\in B-\mathfrak q$, $y$ est intérieur à
$f(\mathrm D(t))$, donc $f(\mathrm D(t))$ contient
$\operatorname{Spec}(\mathcal O_{Y,y})$~; d'où
$\operatorname{Spec}(\mathcal O_{Y,y})
\subset f(\operatorname{Spec}(\mathcal O_{X,x}))$, ce qui établit
notre assertion. Enfin, lorsque $f$ est localement de présentation finie, b)
implique a) en vertu de (1.10.2) appliqué à la restriction $U\to Y$ de $f$ à
un voisinage ouvert arbitraire $U$ de $x$.
\end{proof}

\begin{corollary}[1.10.4]
\label{IV.1.10.4-fr}
Soit $f:X\to Y$ un morphisme. Considérons les conditions suivantes~:
\begin{enumerate}
  \item[a)] $f$ est ouvert.
  \item[b)] Pour tout $x\in X$ et toute générisation $y'$ de $y=f(x)$, il
  existe une générisation $x'$ de $x$ telle que $f(x')=y'$.
  \item[b')] Pour tout $x\in X$, l'image par $f$ de
  $\operatorname{Spec}(\mathcal O_{X,x})$ est
  $\operatorname{Spec}(\mathcal O_{Y,f(x)})$.
  \item[c)] Pour toute partie fermée irréductible $Y'$ de $Y$, toute
  composante irréductible de $f^{-1}(Y')$ domine $Y'$.
\end{enumerate}
Alors on a les implications
$\mathrm{a)}\Rightarrow\mathrm{b)}\Leftrightarrow\mathrm{b')}
\Leftrightarrow\mathrm{c)}$. Si en outre $f$ est localement de présentation
finie, les quatre conditions sont équivalentes.
\end{corollary}

\begin{proof}
Dire que $f$ est ouvert signifie que $f$ est ouvert en tout point $x\in X$,
donc les implications
$\mathrm{a)}\Rightarrow\mathrm{b)}\Leftrightarrow\mathrm{b')}$ résultent des
implications analogues dans (1.10.3), ainsi que l'implication
$\mathrm{b)}\Rightarrow\mathrm{a)}$ lorsque $f$ est localement de présentation
finie. L'implication $\mathrm{c)}\Rightarrow\mathrm{b)}$ résulte aussi de
l'implication analogue dans (1.10.3)~; prouvons enfin que b) entraîne c). En
effet, soit $x'$ le point générique d'une composante
irréductible de $f^{-1}(Y')$, et montrons que $y'=f(x')$ est le point générique
de $Y'$. Soit $y''$ une générisation de $y'$~; il existe par hypothèse une
générisation $x''$ de $x'$ telle que $f(x'')=y''$, et comme
$x''\in f^{-1}(Y')$, on a nécessairement $x''=x'$, donc $y''=y'$, ce qui
achève la démonstration.
\end{proof}

\begin{center}
\emph{(A suivre.)}
\end{center}

\clearpage
